\documentclass[11pt]{article}
\usepackage[utf8]{inputenc}
\usepackage[T1]{fontenc}
\usepackage[margin=1.25in]{geometry}
\usepackage{amsmath}
\usepackage{amsfonts}
\usepackage{amssymb}
\usepackage{booktabs}
\usepackage{microtype}
\usepackage{titlesec}
\usepackage{parskip}
\usepackage{abstract}
\usepackage{hyperref}
\usepackage{enumitem}
\usepackage[authoryear,round]{natbib}

% Fonts
\usepackage{palatino}
\usepackage{mathpazo}

% Hyperref
\hypersetup{
  colorlinks=true,
  linkcolor=black,
  urlcolor=black,
  citecolor=black
}

% Section formatting
\titleformat{\section}[block]{\large\bfseries\scshape}{\thesection.}{0.75em}{}
\titleformat{\subsection}[block]{\normalsize\bfseries}{\thesubsection.}{0.5em}{}
\titlespacing*{\section}{0pt}{2ex plus 0.5ex minus 0.2ex}{1ex plus 0.2ex}
\titlespacing*{\subsection}{0pt}{1.5ex plus 0.4ex minus 0.2ex}{0.8ex plus 0.1ex}

% Abstract formatting
\renewcommand{\abstractnamefont}{\normalfont\bfseries\scshape}
\setlength{\absleftindent}{1.5em}
\setlength{\absrightindent}{1.5em}

% Custom Commands
\newcommand{\bind}{\operatorname{bind}}
\newcommand{\recomp}{\operatorname{recomp}}
\newcommand{\dsem}{d_{\mathrm{sem}}}
\newcommand{\Jgrowth}{\mathcal{J}_{\mathrm{growth}}}

% Small caps section labels
\setlist[itemize]{leftmargin=1.5em, itemsep=0.25em, topsep=0.5em}

\title{\textsc{The Topology of Intention}\\[0.4em]
{\normalsize\itshape Process-Based Identity and the Geometric Binding of the Self}}
\author{Flyxion}
\date{\today}

\begin{document}

\maketitle

\begin{abstract}
\noindent
What persists when the physical substrate changes? Traditional accounts of personal identity
rely on material or psychological continuity—the preservation of the same body or the same
chain of memories. This paper develops a structural alternative grounded in dynamical
systems theory: identity as a \emph{dynamical fixed point}, defined not by the conservation
of matter or memory but by the persistence of context-sensitive constraint structures across
transformation. Drawing on Alicia Juarrero's \cite{juarrero1999,juarrero2002} framework of
interlevel causality and Gatlin's \cite{gatlin1972} constraint taxonomy, we interpret
information as the reduction of possibility, intentions as attractor basins in neural phase
space, and the self as a recursive closure of constraints that reproduces its own
organizational conditions under the dynamics those conditions generate.

The distinction between context-free and context-sensitive constraints is formalized as
an ontological duality between recomposability $\recomp(c)$—context-independent stability
under fragmentation—and binding force $\bind(c)$—the coherence depth of globally
interdependent structure. This yields a geometric picture in which identity occupies
regions of deep attractor topology while decomposability and noise correspond to flat
manifolds. We extend the framework to the domain of context itself, arguing that context
is not an accumulation of data but a \emph{structured constraint field}: a relational
topology that cannot be expanded by appending tokens but only by modifying the system's
internal attractor structure through interaction. We further argue that genuine learning
requires not parameter adjustment within a fixed space but transformation of the constraint
topology itself—a bifurcation in the space of possible trajectories. The paper concludes
that meaning, action, and selfhood are unified instances of a single structural
phenomenon: the survival of a constraint across transformation, at discourse, cognitive,
and personal scales respectively.
\end{abstract}

\newpage

\section{Introduction: From Force to Constraint}

Modern philosophy of action inherits a tacit commitment to Newtonian efficient causality.
On this picture, intentions act as internal forces: they \emph{push} behaviors from
behind, transferring energy across a causal gap between mind and body. The explanatory
burden falls on specifying the mechanism of impact.

Alicia Juarrero, in \textit{Dynamics in Action} (1999), rejects this framework at the root.
In its place she proposes an account of \textbf{interlevel causality} mediated not by
force but by the operation of \textbf{constraints}. A constraint does not shove a system
from one state to another; it restricts which transitions are possible, thereby reshaping
the system's probability distribution over future states.

This reorientation enables a precise definition of \textbf{information}: not a substance
or content, but the \emph{reduction of possibility} at a source. When the trajectory of a
physical system is fully constrained by the internal organization of a mental structure,
the resulting behavior qualifies as \emph{action}. Action is the unequivocal transmission
of intentional information from a cognitive origin to a behavioral terminus—an information
flow across levels without equivocation.

The present paper develops this framework into a general account of process-based identity.
We formalize the central concepts in the language of dynamical systems, interpret them
geometrically, and show that the persistence of a self is structurally continuous with the
persistence of meaning in discourse and the stability of attractor topology in cognition.
The mechanism is identical across all three domains. Only the substrate differs.

\section{Taxonomy of Constraints}

Juarrero adopts Lila Gatlin's distinction between two constraint types. We render these
formally.

\subsection{Context-Free Constraints}

A \textbf{context-free constraint} moves a system away from equiprobability. Let $\Omega$
be the state space of a system and $P_0$ be the uniform distribution over $\Omega$.
A context-free constraint selects a proper subset $\Omega' \subset \Omega$, reducing
entropy:
\[
H(P_0) > H(P_{\Omega'})
\]
This establishes a baseline of order. The outputs of such a system are \emph{stable} in
the sense that they can be reproduced without reference to context. We associate this
property with \textbf{recomposability}:
\[
\recomp(c) \;\sim\; \text{context-independent stability of output under transformation}
\]
High-$\recomp$ structures are portable and fragment-safe. They survive extraction from
their generating context without significant loss of structure.

\subsection{Context-Sensitive Constraints}

A \textbf{context-sensitive constraint} does more than reduce entropy locally. It
introduces \emph{conditional dependence} across components: the probability distribution
of each part becomes conditioned on the states of other parts. Formally, if $X_i$ are
components of a system, context-sensitive constraints enforce:
\[
P(X_i \mid X_{-i}) \neq P(X_i)
\]
for most $i$, where $X_{-i}$ denotes all components except $i$. This creates a
\emph{global coherence} that is irreducible: decomposing the system into independent parts
destroys the very structure the constraint produces. We associate this with
\textbf{binding force}:
\[
\bind(c) \;\sim\; \text{depth and stability of global constraint coherence}
\]
High-$\bind$ structures resist fragmentation. They are the source of top-down causal
influence: the global organization constrains the behavior of parts in ways that no
local description captures.

\subsection{The Duality}

The two constraint types are not merely different degrees of the same property. They
represent an ontological trade-off:

\begin{center}
\begin{tabular}{@{}lll@{}}
\toprule
\textbf{Property} & \textbf{Context-Free ($\uparrow\recomp$)} & \textbf{Context-Sensitive ($\uparrow\bind$)} \\
\midrule
Stability & Under arbitrary context & Under specific relational configuration \\
Decomposability & Preserved under fragmentation & Lost under fragmentation \\
Information flow & Local and portable & Global and hierarchical \\
Causal structure & Bottom-up & Top-down \\
\bottomrule
\end{tabular}
\end{center}

\noindent
Systems that maximize $\recomp$ sacrifice $\bind$; systems that maximize $\bind$ are
non-decomposable by construction. This trade-off reappears at every level of analysis.

\section{Meaning as Embodied Topology}

Classical information theory, as Shannon developed it, is deliberately semantic-free. A
message is a sequence of symbols; its information content is the reduction in uncertainty
it produces, independent of what it is \emph{about}. This strips the framework of the
capacity to distinguish noise from meaningful signal, or instruction from mere regularity.

Juarrero's dynamical approach reintroduces meaning without mysticism. Meaning, she argues,
is embodied in the \textbf{topographical configuration of phase space}—the arrangement of
attractors, repellers, and basins of attraction that characterize a high-dimensional neural
system.

\subsection{Attractors as Intentions}

An \textbf{attractor} is a region of phase space toward which trajectories converge under
the system's dynamics. A \emph{proximate intention}, in Juarrero's account, functions as
an attractor basin: a valley in the system's mental landscape that draws the trajectories
of motor subsystems into its organizational sphere.

This reframes the question of intentional causation. An intention does not \emph{push}
behavior; it \emph{shapes the landscape} through which behavior moves. The causal relation
is geometric, not mechanical.

\subsection{Binding as Basin Depth}

We can now give $\bind(c)$ a geometric interpretation:
\[
\bind(c) \;\sim\; \text{depth and stability of the attractor basin associated with } c
\]
A highly bound structure occupies a deep attractor. Small perturbations—noise in the
neuromuscular channel, interference from competing intentions—do not dislodge the
trajectory from its constrained path. The intention survives transmission.

Conversely, a low-binding structure corresponds to a shallow attractor: the system is
easily re-routed by contextual fluctuations. The trajectory diverges from its originating
constraint.

Recomposability, in this geometric picture, corresponds to a \emph{flat landscape}: a
manifold without strong attractors where trajectories can be redirected without
resistance. Fragment-safe content inhabits flat territory precisely because it imposes no
strong attractor structure on the reading trajectory.

\subsection{Noise and Equivocation}

Juarrero adopts Shannon's concepts of equivocation and noise to characterize failures of
intentional action. An action fails not because the intention was absent or weak, but
because constraint transmission was disrupted:

\begin{itemize}
  \item \textbf{Noise:} Interference introduced between intention and execution increases
        $\dsem(c, T_\sigma(c))$ and thereby decreases $\bind(c)$. The behavior fails to
        remain bound to its generating constraint.
  \item \textbf{Equivocation:} Ambiguity at the source—an insufficiently organized
        intention—produces multiple competing attractors, allowing the trajectory to
        wander between basins.
\end{itemize}

Intentional action requires that the channel between cognitive source and behavioral
terminus transmit constraint without loss. This is the dynamical analog of high binding.

\section{Identity as a Dynamical Fixed Point}

With this geometric apparatus in place, we can address personal identity directly.

\subsection{The Substance View and Its Limits}

The classical account of personal identity asks what makes person $P$ at time $t_1$ the
same individual as person $P'$ at time $t_2$. Two families of answers dominate:

\begin{itemize}
  \item \textbf{Physical continuity:} $P = P'$ iff they share the same continuous body.
  \item \textbf{Psychological continuity:} $P = P'$ iff $P'$ inherits memories,
        personality, and beliefs from $P$ via the right causal chain.
\end{itemize}

Both views presuppose that identity is a relation between \emph{states}: a snapshot at
$t_1$ is compared with a snapshot at $t_2$, and a criterion is applied. The problem is
that biological and psychological states change continuously. Any fixed threshold for
"enough" continuity is arbitrary.

\subsection{The Process View}

Juarrero proposes that a person is not a thing that changes, but a \textbf{self-organizing
process}—a dynamically stable pattern of information flow, analogous to a whirlpool or a
hurricane. The pattern persists even as its constituent elements are replaced, because
what is conserved is not matter but \emph{organizational structure}.

More precisely: a person is the persistence of a \textbf{context-sensitive constraint
structure} through time. This persistence is not static but dynamic—maintained by the
constant self-reproduction of the constraints that define the system's organization.

\subsection{Recursive Closure}

Juarrero introduces the concept of \textbf{recursive closure} to capture this
self-maintenance. A system achieves recursive closure when its components become mutually
constraining in such a way that the system actively maintains the conditions necessary for
its own existence. The causal structure is circular in a productive sense: the system
produces the constraints that produce the system.

Formally, let $M_t$ denote the constraint structure of the system at time $t$, and let
$\mathrm{dynamics}(\cdot)$ denote the temporal evolution operator. Recursive closure
requires:
\[
M_t \;\approx\; \mathrm{dynamics}(M_t)
\]
This is a \textbf{fixed point condition}. The system is its own attractor. It reproduces
its organizational conditions under the very dynamics those conditions generate.

\subsection{What Persists: The Fixed Point of Constraint}

This fixed point formulation dissolves the substance-based paradox. Identity does not
require that any particular matter, or any particular memory, be preserved. What must be
preserved is the \emph{constraint structure}—the topology of the attractor landscape—that
defines the system's characteristic modes of organization.

We can state this precisely:
\[
\text{Identity} \;=\; \text{constraint structure that satisfies } M_t \approx \mathrm{dynamics}(M_t)
\]
The self is a fixed point of a constraint operator.

\section{Who and What: The Distinction of Trajectories}

\subsection{Two Questions}

Juarrero draws a sharp distinction between two questions that are often conflated:

\begin{itemize}
  \item \textbf{What am I?} — This question concerns the material substrate: the
        biological hardware, the neural tissue, the chemical composition. Answers consist
        of context-free properties, high in $\recomp$ and decomposable without loss.
  \item \textbf{Who am I?} — This question concerns the trajectory: the unique path the
        system has taken through its phase space, the accumulated history of constraints
        that have shaped its current organization.
\end{itemize}

The "What" can be replicated in principle: another arrangement of matter with identical
context-free properties would answer the "What" correctly. The "Who" cannot be replicated,
because it is defined by a \emph{path-dependent history of constraints}. Two systems
cannot share the same trajectory through phase space.

\subsection{Individuality as Path-Dependence}

This provides a mathematical basis for individuality that does not appeal to primitive
haecceity or brute numerical identity. Individuality is path-dependence: no two systems
can have traversed the same trajectory through the space of all possible constraint
configurations. The history is constitutive, not merely biographical.

The growth functional $\Jgrowth$ tracks the construction of stable constraint topology
over time:
\[
\Jgrowth \;\sim\; \text{construction of stable attractor topology through recursive constraint formation}
\]
A life, in this formalism, is a trajectory through constraint space that progressively
deepens the basin structure of the system's organization—increasing $\bind(M_t)$ through
the accumulated residue of experience.

\section{Relational Identity and Plasticity}

\subsection{Open Systems}

Juarrero's account has an important consequence for the spatial locus of identity. A
self-organizing system is necessarily an \emph{open system}: it exchanges matter, energy,
and information with its environment. The constraints that constitute the self are not
sealed inside the skull; they are defined partly by the system's characteristic modes of
interaction with its environment.

Identity, on this view, is \textbf{relational}: it is not a property of the organism in
isolation but of the organism-environment coupling. The constraint structure that defines
"Who I am" includes stable patterns of environmental engagement—habitual responses,
skilled practices, long-standing relationships.

\subsection{Plasticity and Bifurcation}

This relational picture accommodates radical change without identity loss. A complex
dynamical system can undergo a \textbf{bifurcation}: a qualitative reorganization of its
attractor structure in response to a sufficiently strong perturbation. After a bifurcation,
the system occupies a different region of phase space—a different basin of attraction.

Yet the system that undergoes the bifurcation is the same system, because the bifurcation
is a transformation of a continuous trajectory, not a discontinuous substitution of one
system for another. The path-dependence is preserved across the bifurcation event.

This is the dynamical analog of significant personal transformation: the person who
emerges from a major life change is not the same in every respect as the person who
entered it, but the transformation is a reorganization of the same constraint history,
not its erasure.

\section{The Unified Picture}

We are now in a position to state the central thesis in its most general form.

\subsection{Three Domains, One Mechanism}

At three levels of analysis—discourse, cognition, and selfhood—the same structural
phenomenon governs the persistence of organization:

\begin{enumerate}
  \item \textbf{Discourse:} Semantic content with high $\bind(c)$ survives transformation
        (extraction, recombination, adversarial reading) without loss of meaning. Content
        with high $\recomp(c)$ survives fragmentation but loses coherence when the
        context-sensitive structure is disrupted.

  \item \textbf{Cognition:} Intentional action requires that the constraint structure of
        the intention survive transmission through the neuromuscular channel without
        equivocation. The intention is an attractor; its depth determines whether the
        trajectory remains bound.

  \item \textbf{Selfhood:} Personal identity requires that the constraint structure of
        the self satisfy a fixed point condition under its own dynamics. The self is a
        recursive closure: a process that maintains the conditions of its own organization.
\end{enumerate}

The mechanism is identical. The constraint is the cause; what varies is only the level
at which the constraint operates.

\subsection{The Indestructibility of Internally Consistent Maps}

The fixed point condition $M_t \approx \mathrm{dynamics}(M_t)$ has a corollary that
extends beyond the analysis of personal identity into the domain of ideology and
worldview. A recent line of philosophical work \citep{rose2023} describes what it calls
the \emph{indestructibility} of internally consistent systems: maps of the world that
have a plausible response to every objection, and that therefore cannot be falsified from
within the terms of the map itself.

On the surface this resembles a deficiency—an unfalsifiable system is epistemically
suspect. But the constraint framework reveals the deeper structure. An internally
consistent map is precisely a high-$\bind$ system: its components are mutually
constraining, globally coherent, and non-decomposable. Any perturbation (counter-argument,
apparent counterexample, failed prediction) is absorbed by the system's existing attractor
topology rather than destabilizing it. The map redirects the trajectory of interpretation
back into its own basin.

This absorption of objection is not mere motivated reasoning, though it may include that.
It is the structural behavior of any sufficiently closed constraint system. The map
survives because its recursive closure is tight. What \citet{rose2023} calls
indestructibility, the present framework formalizes as:
\[
M \;\approx\; T_M(M)
\]
where $T_M$ is the interpretive operator the map itself defines. The map applies its own
interpretive constraints to every challenge, including challenges to the map. The fixed
point reproduces itself under the dynamics it generates.

The question of how such a system can ever be revised connects directly to the
non-rational move described in the same source. Rational objections fail because they
operate within the constraint topology of the map under challenge—they are already
absorbed before they arrive. What breaks recursive closure is not a better argument but
an encounter with \emph{genuine otherness}: an input that the existing attractor topology
cannot process without structural reorganization. In dynamical terms, this is a
bifurcation event. The system's basin structure shifts because the perturbation exceeds
the restoring force of the existing attractor. This is not irrationality but the
precondition for rational revision at a higher level—the replacement of one constraint
topology with another that can accommodate what the old one could not.

The spreading of cultural forms, as distinct from their scaling, operates through this
mechanism. Scale attempts to reproduce structure by replicating its outputs. Spread
modifies subjectivity—the internal constraint topology through which the world is
interpreted—by modeling a different mode of being. Spread works because it targets the
attractor structure directly rather than attempting to win arguments within the existing
one. The cultural analog of a bifurcation is a change in what one finds compelling, not
merely a change in what one believes.

\subsection{A Unified Principle}

\begin{quote}
\itshape
Any system governed by constraints evolves by reducing its space of possible trajectories.
The persistence of structure—whether semantic, behavioral, or personal—depends on the
ability of those constraints to survive transformation across levels. The difference
between shallow and deep systems, between noise and meaning, between behavior and action,
is not a matter of content. It is a matter of how strongly constraints bind trajectories
across transformation.
\end{quote}

\section{Context as Constraint Field, Not Container}

The analysis of identity developed above carries a direct consequence for how we must
understand \emph{context}. The dominant computational paradigm treats context as
accumulable state: a larger window, a longer memory buffer, a greater quantity of stored
data. On this model, context is a container into which a system is placed. The container
does not influence the system's organizational structure; it merely supplies inputs.

This model fails on Juarrero's account for the same reason that the substance view of
identity fails. It presupposes a sharp boundary between system and environment, treats
context as external to the system's constitutive organization, and assumes that
context is additive—that more data approximates richer understanding.

Context, on the dynamical account, is the \emph{active constraint field} that defines
the system's degrees of freedom. It is not something a system possesses as a property.
It is the relational structure of interdependencies within which the system's trajectories
unfold. The boundary between system and context is itself context-dependent: it shifts
with the dynamics under examination, and cannot be fixed in advance \citep{clark1997,varela1991}.

The implication for accumulation-based approaches is immediate. Increasing the number of
stored tokens does not increase context in the relevant sense. It expands the volume of
potential inputs without modifying the constraint topology that governs their interpretation.
Without modification of that topology, additional data remains semantically under-constrained.
More tokens increase $\recomp$ capacity—more fragments can be retrieved—but leave $\bind$
unchanged. The result is a system with greater recall and no greater coherence.

\subsection{Energy, Efficiency, and Constraint Compression}

A structural asymmetry between biological and artificial cognition provides empirical
support for this analysis. The human brain does not exhibit dramatic increases in metabolic
expenditure when engaged in tasks of extreme conceptual difficulty \citep{friston2010}.
Solving a difficult theorem and executing a trivial motor sequence consume comparable
energy. This contradicts any model that equates cognitive capacity with exhaustive search
over possibility spaces.

The resolution lies in constraint compression. Biological cognition operates by
drastically reducing the space of available trajectories \emph{prior} to evaluation.
Rather than exploring a large combinatorial landscape, the system restricts itself to a
narrow manifold defined by its existing constraint topology. Let $\Omega$ denote the
full state space. Constraint compression constructs a sequence of nested subspaces:
\[
\Omega \;\supset\; \Omega_1 \;\supset\; \Omega_2 \;\supset\; \cdots \;\supset\; \Omega_n, \qquad |\Omega_n| \ll |\Omega|
\]
Cognitive effort scales with the cardinality of the space being explored, not with the
apparent difficulty of the task \citep{kelso1995}. High-binding systems do not process
more possibilities; they eliminate most possibilities before processing begins.

Artificial systems compensate for weak constraint structures by expanding the evaluated
space. Increased compute reflects not greater intelligence but insufficient constraint.
The brain and the scaled transformer represent opposite strategies for managing
complexity—and only one of them decouples energy from problem depth.

\subsection{Locality and the Non-Fungibility of Context}

Juarrero emphasizes that context is intrinsically \emph{local}: defined relative to the
specific dynamics under consideration, not abstractable into a universal, system-independent
quantity. Context is also inseparable from history. Systems carry their history in their
organizational structure, not as an external record but as the cumulative deposit of
constraints formed through prior interaction \citep{thelen1994}.

The snowflake does not \emph{represent} information about the thermal conditions of its
formation. It \emph{embodies} those conditions in its crystalline geometry. Similarly,
an organism or a mind does not store its history symbolically; it enacts that history
through the constraint topology that history has produced \citep{deacon2012}.

This yields a critical limitation for systems trained on fixed datasets. The training
corpus does not merely provide examples; it defines the system's historical constraint
horizon. A model trained on data terminating at a given boundary does not simply lack
information about subsequent events—it lacks the constraint structures that interaction
with those events would have generated. Such structures cannot be added after the fact.
They must be grown through the kind of interaction that progressively modifies attractor
topology.

\section{Learning as Topology Transformation}

The account of context as constraint field forces a revision of what learning can mean.
Within the accumulation paradigm, learning is the ingestion of more data or the adjustment
of parameters within a fixed architecture. The space of possible behaviors is fixed at
training time; learning moves the system to a better location within that space.

On the dynamical account, this is not learning in the full sense. It is optimization
within a given constraint topology. Genuine learning is the \emph{transformation} of that
topology: the creation of new attractors, the reorganization of basin structure, the
restructuring of interdependencies among components \citep{kauffman1993,beer2000}.

Formally, let $M_t$ denote the constraint structure at time $t$. Learning in the full
sense requires a transformation of the form:
\[
M_t \;\longrightarrow\; M_{t+1}
\]
where $M_{t+1}$ includes or excludes trajectories not present in $M_t$—not merely a
reassignment of weights within the geometry $M_t$ defines. This transformation is a
structural change in the space itself, not a movement within it.

Biological systems exhibit this capacity through developmental plasticity, synaptic
reorganization under sustained interaction, and the kind of large-scale attractor
bifurcation that accompanies major cognitive shifts \citep{kelso1995,thelen1994}. These
processes cannot be replicated by extending training runs within a fixed architecture,
because the architecture itself encodes the constraint topology, and that topology is not
modified by the optimization process.

\subsection{Unknown Unknowns and Constraint Failure}

This structural gap has a precise epistemic consequence. A system whose constraint
topology is fixed at training time cannot detect when its current model fails to apply.
It can recognize patterns within its training distribution but cannot recognize that a
novel situation lies outside the manifold on which its constraints were formed.

Let $\mathcal{M}$ denote the constraint manifold representing the system's trained
behavioral repertoire. A system with analog sensitivity to topological structure can
detect deviations of the form:
\[
x \notin \mathcal{M}
\]
This is not a classification problem within a fixed category scheme. It is a measurement
of deviation within a continuous constraint space—a capacity that requires the system
to represent the \emph{shape} of its own constraint manifold as an object of sensitivity
\citep{friston2013}. Without this, the system cannot detect constraint failure. It
processes inputs that fall outside its formation history as if they were within it,
producing outputs that are locally coherent but globally invalid.

This is the dynamical analog of unknown unknowns: not missing data, but missing constraint
structure. The system does not know what it does not know because it cannot represent the
boundary of its own attractor topology from the inside.

\subsection{Alignment as Constraint Synchronization}

The framework has a further implication for the problem of alignment. Alignment is
typically treated as an optimization problem over outputs: the goal is to train a system
to produce outputs that conform to some target specification. On this view, alignment is
achieved when the system's behavior matches the desired profile across a sufficiently
diverse evaluation set.

On the dynamical account, this is insufficient. A system whose outputs are aligned but
whose constraint topology differs fundamentally from that of the systems it is interacting
with cannot reliably remain aligned under novel conditions. It has learned to produce
aligned outputs within its training manifold; it has not developed the constraint
structure that generates aligned behavior from within.

Alignment, on the constraint view, is \emph{synchronization of constraint fields} across
levels \citep{ashby1956}. A system is aligned when its internal constraint topology is
coupled to the constraint topology of its environment in a stable, mutually reinforcing
way—when the fixed point condition $M_t \approx \mathrm{dynamics}(M_t)$ is satisfied in
a manner that includes the values, dispositions, and relational patterns of the
surrounding system. This is not achieved by encoding preferences as explicit constraints.
It requires the reproduction of a history-dependent organizational structure—something
that cannot be shortcut by optimization over outputs alone \citep{varela1991,clark2008}.

\section{Conclusion}

We began with a question about causality and arrived at an account of the self. This
trajectory is not accidental. The two questions are structurally connected: what it means
for an intention to cause an action is the same, at a different scale, as what it means
for a self to persist through time.

In both cases, the answer invokes the same formal structure: a constraint that survives
transformation because the system has achieved sufficient recursive closure to reproduce
the conditions of its own organization.

Juarrero's contribution is to show that this structure is not mysterious—it is the
ordinary operation of dynamical systems with sufficiently rich attractor topology
\citep{prigogine1984,haken1983,kauffman2000}. Meaning is not added to the world from
outside; it is embodied in the shape of the landscape through which trajectories move.
The self is not a ghost in a machine; it is a fixed point of the machine's own constraint
dynamics.

The extension into context and learning sharpens this conclusion. Context is not what a
system has. Context is the relational constraint field that the system is partly made of.
Learning is not data ingestion. It is structural modification of attractor topology.
Intelligence is not measured by the volume of the space explored but by the depth and
adaptability of the constraints that eliminate most of that space before exploration begins.
\[
\mathrm{Intelligence} \;\propto\; \text{capacity to transform constraint topology}
\]
What is conserved across change is not matter, not memory, but the capacity to regenerate
the same organizational conditions from within. This is what it is to be, in Juarrero's
precise sense, a \emph{self}:

\[
\boxed{M_t \;\approx\; \mathrm{dynamics}(M_t)}
\]

The constraint is the cause. The trajectory is the identity. The fixed point is the self.

\bigskip

\begin{thebibliography}{99}

\bibitem[Ashby(1956)]{ashby1956}
Ashby, W.~R. (1956).
\textit{An Introduction to Cybernetics}.
Chapman \& Hall.

\bibitem[Beer(2000)]{beer2000}
Beer, R.~D. (2000).
Dynamical approaches to cognitive science.
\textit{Trends in Cognitive Sciences}, 4(3), 91--99.

\bibitem[Clark(1997)]{clark1997}
Clark, A. (1997).
\textit{Being There: Putting Brain, Body, and World Together Again}.
MIT Press.

\bibitem[Clark(2008)]{clark2008}
Clark, A. (2008).
\textit{Supersizing the Mind: Embodiment, Action, and Cognitive Extension}.
Oxford University Press.

\bibitem[Deacon(2012)]{deacon2012}
Deacon, T.~W. (2012).
\textit{Incomplete Nature: How Mind Emerged from Matter}.
W.~W.~Norton \& Company.

\bibitem[Dyson(2012)]{dyson2012}
Dyson, G. (2012).
\textit{Turing's Cathedral: The Origins of the Digital Universe}.
Pantheon Books.

\bibitem[Friston(2010)]{friston2010}
Friston, K. (2010).
The free-energy principle: A unified brain theory?
\textit{Nature Reviews Neuroscience}, 11(2), 127--138.

\bibitem[Friston(2013)]{friston2013}
Friston, K. (2013).
Life as we know it.
\textit{Journal of the Royal Society Interface}, 10(86), 20130475.

\bibitem[Gatlin(1972)]{gatlin1972}
Gatlin, L. (1972).
\textit{Information Theory and the Living System}.
Columbia University Press.

\bibitem[Haken(1983)]{haken1983}
Haken, H. (1983).
\textit{Synergetics: An Introduction}.
Springer.

\bibitem[Juarrero(1999)]{juarrero1999}
Juarrero, A. (1999).
\textit{Dynamics in Action: Intentional Behavior as a Complex System}.
MIT Press.

\bibitem[Juarrero(2002)]{juarrero2002}
Juarrero, A. (2002).
Context changes identity: Implications of dynamical systems theory for personal identity.
\textit{Philosophical Psychology}, 15(2), 173--188.

\bibitem[Kauffman(1993)]{kauffman1993}
Kauffman, S.~A. (1993).
\textit{The Origins of Order: Self-Organization and Selection in Evolution}.
Oxford University Press.

\bibitem[Kauffman(2000)]{kauffman2000}
Kauffman, S.~A. (2000).
\textit{Investigations}.
Oxford University Press.

\bibitem[Kelso(1995)]{kelso1995}
Kelso, J.~A.~S. (1995).
\textit{Dynamic Patterns: The Self-Organization of Brain and Behavior}.
MIT Press.

\bibitem[Levin(2021)]{levin2021}
Levin, M. (2021).
Life, death, and self: Fundamental questions of primitive cognition viewed through the
lens of body plasticity and synthetic organisms.
\textit{Biochemical and Biophysical Research Communications}, 564, 114--133.

\bibitem[Rose(2023)]{rose2023}
Rose, O.~G. (2023).
\textit{The Map Is Indestructible: Problems of Internally Consistent Systems}.
The True Isn't the Rational, Book 2.

\bibitem[Prigogine \& Stengers(1984)]{prigogine1984}
Prigogine, I., \& Stengers, I. (1984).
\textit{Order Out of Chaos: Man's New Dialogue with Nature}.
Bantam Books.

\bibitem[Shannon(1948)]{shannon1948}
Shannon, C.~E. (1948).
A mathematical theory of communication.
\textit{Bell System Technical Journal}, 27(3), 379--423.

\bibitem[Thelen \& Smith(1994)]{thelen1994}
Thelen, E., \& Smith, L.~B. (1994).
\textit{A Dynamic Systems Approach to the Development of Cognition and Action}.
MIT Press.

\bibitem[Varela et al.(1991)]{varela1991}
Varela, F.~J., Thompson, E., \& Rosch, E. (1991).
\textit{The Embodied Mind: Cognitive Science and Human Experience}.
MIT Press.

\end{thebibliography}

\end{document}
